\documentclass[12pt]{ctexart}
\author{Yan}
\usepackage{amsmath,amssymb}
\usepackage{geometry}
\usepackage{setspace}
\usepackage{multirow}
\usepackage{array}
\usepackage{makecell}
\geometry{a4paper, margin=1in}

\geometry{a4paper,margin=2.5cm}
\setstretch{1.25}

\title{拍拍 Clapclap 2.0 规则书}
\date{2026.3.20}

\begin{document}

\maketitle
\tableofcontents
\newpage

%------------------------------------------------
\section{游戏介绍}

拍拍（Clapclap）是一款多人参与的同步出招战斗游戏。所有玩家围成一圈，通过击掌与手势的组合进行回合制对抗。
游戏中的所有玩家在同一时刻完成出手，并在统一的结算阶段中处理所有效果。
玩家需要根据自己拥有的资源、场上其他玩家可能做出的选择，以及多人结算时的先后关系来决定本回合的手势。

在游戏中，玩家会不断在资源积累、攻击、防御与锦囊反制之间进行选择。
部分手势偏向发育，部分手势用于直接攻击，部分手势可以抵挡伤害，还有部分手势能够针对特定攻击方式进行反制或改变局势。
游戏的进行方式并不是传统的轮流行动，而是所有人先同时出手，再统一进入结算，
因此每回合的结果不仅取决于自己出了什么，也取决于其他玩家出了什么，以及这些手势在多人规则中的结算顺序。

本规则书介绍的是拍拍 2.0 多人版规则。与双人版相比，多人版最大的变化不在于手势本身，而在于结算阶段的处理方式。
多人版中引入了“未操作对象”“已操作对象”“速度”“三连”等概念，用于决定谁先结算、谁可以被选为目标、哪些手势会互相抵消，
以及在复杂局面中如何保证整个回合能够被清晰地处理完成。

游戏的目标是：在不断进行的回合中存活下来，直到场上只剩最后一名玩家，该玩家获得胜利。

%------------------------------------------------
\section{资源介绍}

游戏中一共涉及六种资源，分别是：生命值、气、盾、火种、电池与镐。

在默认规则下，除镐外，其余资源均没有上限。

\subsection{生命值}

生命值代表玩家当前的生存状态。

在一个回合结束时，若一名玩家的生命值降至 0 及以下，该玩家视为死亡。
并在下一回合开始前视为已经出局，不再参与后续回合。

若场上最后只剩一名存活玩家，则游戏结束。

\subsection{气}

气是本游戏最基础的资源之一。

它可以通过资源手势积累，也可以被消耗来发动若干攻击手势、防御手势与锦囊手势。
很多重要的战斗手势都与气有关，因此气是游戏中极为常见的一类资源。

\subsection{盾}

盾是另一种基础资源。

它同样可以通过资源手势积累，并主要用于发动盾系攻击。
与气相比，盾本身的用途更集中在某一类攻击路线中，但它也会衍生出火种与电池两种中间资源。

\subsection{火种}

火种是由部分盾系攻击衍生出来的资源。

它不能通过普通资源手势直接获得，而是要通过特定攻击手势Fire产生，烈焰可以直接消耗火种发动。

\subsection{电池}

电池与火种类似，也是一种由部分盾系攻击衍生出来的资源。

它不能通过普通资源手势直接获得，而是要通过特定攻击手势闪电产生，Shining可以直接消耗电池发动。

\subsection{镐}

镐与一般资源不同，它不只是单纯的储存型资源，还会直接影响玩家的生命状态，并带有额外风险，是游戏中的特殊资源。

当玩家受到伤害时，如果身上有镐，则会优先消耗镐来抵挡伤害，每消耗一个镐可以抵挡掉 1 点伤害。
若玩家在获得镐的时候生命值为 0 或以下，则会直接恢复 1 点生命值，而不是获得一个镐。

若一名玩家持有镐的数量超过了上限，则会立即触发“爆镐”，该玩家直接死亡，该上限一般为 1。

%------------------------------------------------
\section{游戏进行}


所有玩家在游戏中需要保持固定的站位与基本姿势：玩家围成一个闭环，每个人都同时与左右两侧玩家发生拍手关系。
为了方便击掌与保持游戏动作统一，每名玩家需要将右手手掌朝下、左手手掌朝上，使自己能够自然地与左右两侧玩家完成拍手动作。

游戏开始时，每位玩家都拥有相同的初始状态。
通常情况下，每位玩家初始拥有 1 点生命值，其余资源均为 0。
若某一局有额外约定，也可以在开局前统一说明，但在默认规则下，所有人都以相同资源开始游戏。

游戏按回合进行。每一回合由三个阶段组成，依次为：击掌阶段、出手阶段、结算阶段。

\subsection{击掌阶段}

在击掌阶段中，所有玩家同时与左右两侧玩家各上下击掌一次。
由于玩家是围成一圈站立或坐下的，因此这一动作会在全体玩家之间形成一个闭环的同步互动。

击掌阶段本身不直接造成伤害，也不消耗资源，它的作用在于形成游戏的固定节奏，并作为每回合出手前的统一动作。

\subsection{出手阶段}

在出手阶段中，所有玩家必须同时做出一个手势，每位玩家每回合只能出一个手势。
在多人游戏中，出手阶段的同步性非常重要，出手时不能延迟，
不能在观察到其他玩家的动作后再临时改手势，也不能先出一个模糊动作再根据他人动作补正。

所有玩家的出手必须尽量保持同步，并且手势必须符合规则书中允许使用的正式手势。
不能延迟或观察他人后再决定，慢出手的玩家被称作“蛤蟆”，出不合规手势的玩家被称作“蟆蛤”，两者都将直接死亡。

\subsection{结算阶段}

结算阶段是多人版规则的核心部分。当所有玩家都完成出手后，游戏进入统一结算，所有结算都围绕本回合已经亮出的手势展开。

玩家们首先需要观察所有人的手势，确认是否有人消耗了过量资源，
消耗过量气/盾的行为称为“爆气”/“爆盾”，消耗过量气/盾的玩家被称作“蚂蚁”，并直接死亡。

之后，玩家们需要根据本回合所有亮出的手势，按照规则书中规定的结算机制来处理每个手势的效果。

%------------------------------------------------
\section{基础判定}

\subsection{攻击力与防御力}

在本游戏中，绝大多数结算都要围绕攻击力和防御力展开。

攻击力表示一个手势在攻击时的强度。
只有攻击手势具有攻击力，通常会有明确的攻击力数值，例如 1、2、3、4、5 或 1.5 等。

防御力表示一个手势在面对攻击时的抵挡能力。
所有攻击手势的防御力等于其攻击力，例如，若某攻击手势的攻击力为 3，则它作为被攻击对象时，其防御力也视为 3。
资源手势与防御手势的防御力由规则单独规定。锦囊手势的防御力为0。

\subsection{基础攻防判定}

本游戏最基础的攻防判定规则是：

\begin{itemize}
    \item 当一方的攻击力大于目标的防御力时，攻击成立；
    \item 当一方的攻击力小于攻击手势的攻击力或资源手势的防御力时，无法选择其为目标进行攻击；
    \item 当一方的攻击力小于防御手势的防御力时，无法造成伤害；
    \item 当双方攻击力相等，且彼此互相攻击时，双方视为“对掉”。
\end{itemize}

攻击成立表示该攻击成功命中并造成该攻击对应的伤害，而无法造成伤害并不总是等于什么都没发生。
因为在多人版中，某些未打穿防御的攻击仍然可能被视为已经完成一次交互，这一点会在后文结合多人规则具体说明。

“对掉”是本游戏中一个需要特别理解的概念。
在多人版里，若两名玩家使用相同攻击力的攻击手势互相攻击，则允许他们互相攻击，但最终效果等价于双方都没有对彼此进行攻击。
尽管如此，这次交互仍然成立，因此双方都会成为已操作对象。

\subsection{攻击后的防御变化}

在多人版中，一旦一名玩家选择对他人发动攻击，那么该玩家自身的防御力会在攻击发出的同时变为 0。
这个规则只在玩家真正选择攻击别人时生效，而不是只要出了攻击手势就自动生效。
可以理解为攻击打出去之后就离开了自身，因此不再享受该攻击手势原本提供的防御力。

这一规则对同速度层中的互相攻击非常重要。假设两名玩家同属一个速度层并且都使用攻击手势，那么如果其中一人选择主动攻击另一人，
他在攻击打出的同时就会失去原本由该攻击手势提供的防御力，因此可能被另一名同速度玩家利用这一点进行反击。
正因如此，多人版中同速度玩家之间往往需要先声明目标意图，再经过协商决定如何结算。

%------------------------------------------------
\section{手势介绍}

游戏中的所有手势可以分为四大类：资源、攻击、防御和锦囊。

\subsection{资源手势}

资源手势主要用于积累气或盾。它们通常不直接造成伤害，但会为后续的攻击、防御与锦囊提供资源基础。

\subsubsection{气}

手势为：双手作鼓掌状。

获得 1 格气，防御力为 0。

气是游戏中最基础的资源之一。大量攻击手势、防御手势和锦囊手势都需要消耗气，因此气的积累在整个对局中十分重要。
由于气的防御力为 0，使用气时本回合几乎不具备抵挡攻击的能力。

\subsubsection{盾}

手势为：大臂紧贴身体，双拳握于胸前，手背朝向前方。

获得 1 格盾，防御力为 1.5。

盾是另一类基础资源。它主要用于发动盾系攻击。
与气相比，盾在积累资源的同时也提供一定防御，因此出盾通常比出气更不容易被低强度攻击直接打穿。

\subsection{攻击手势}

攻击手势用于主动对他人造成伤害，或在多人规则下对特定目标形成有效威胁。
攻击手势具有明确的攻击力，攻击手势本身的防御力等于其攻击力。

\subsubsection{gi}

手势为：大臂紧贴身体，小臂与大臂垂直，双手握拳向前，手背向下。

消耗 1 格气，攻击力为 1，可造成 1 点伤害。

gi 是最基础的气系攻击手势。它消耗低，使用频率高，且具有多项特殊规则。

\subsubsection{破}

手势为：双手手掌朝前，向前推出。

消耗 2 格气，攻击力为 2，造成 1 点伤害。

破比 gi 更强，但也更容易被针对。

\subsubsection{冷锋}

手势为：右手三根手指作手枪状。

消耗 3 格气，攻击力为 3，造成 1 点伤害。

冷锋属于中强度的稳定攻击，攻击力高于破，能打穿更多低中强度目标，同时也会受到更高攻击与防御的制约。

\subsubsection{如来}

手势为：右手手掌向左，除食指弯曲之外其它手指竖直向上。

消耗 5 格气，攻击力为 4，造成 2 点伤害。

如来是高强度的气系攻击，能打穿大部分防御，但同时也需要较多的气作为发动条件，因此在使用时需要考虑资源积累与时机选择。

\subsubsection{黑洞}

手势为：左手小臂平行置于右手小臂上方，双手指尖置于肘前区形成朝前方的矩形。

消耗 8 格气，攻击力为 5，造成 3 点伤害。

黑洞是气系攻击中攻击力最高的一种。
正常情况下，黑洞命中时造成 3 点伤害。
黑洞还具有拆分能力，可以拆成三个小黑洞分别结算。拆分后的每个小黑洞各造成 1 点伤害，但攻击力仍视为 5。

\subsubsection{Fire}

手势为：双手交叉于腹前，手掌向上手指自然弯曲。

消耗 2 格盾，攻击力为 1.5，造成 1 点伤害，并获得 1 个火种。

Fire 是盾系中最基础的攻击方式，它攻击力不高，但会产生火种为后续烈焰等手势提供资源支持。

\subsubsection{闪电}

手势为：双手置于身前，手心向下，伸出食指和中指并拢，用右手双指第二指关节敲击左手双指第二指关节。

消耗 3 格盾，攻击力为 2，造成 1 点伤害，并获得 1 个电池。

闪电是盾系中的常用攻击方式，可以获得电池为后续Shining等手势提供资源支持。
闪电还可以由 Shining 拆分而来，因此在多人结算中地位较特殊。

\subsubsection{烈焰}

手势为：双手交叉于腹前，手掌向下，手指自然弯曲。

消耗 4 格盾或 2 个火种，攻击力为 3，造成 1 点伤害。

烈焰是 Fire 路线的上位攻击，它比闪电更强，能够打穿更多中强度目标。

\subsubsection{Shining}

手势为：双手各三根手指作手枪状置于额头前上方区域，手背向外，用右手双指第二指关节敲击左手双指第二指关节。

消耗 6 格盾或 2 个电池，攻击力为 4，造成 2 点伤害。

Shining 是盾系中的高阶攻击，可以拆分为两个闪电分别结算。

\subsection{防御手势}

防御手势用于抵挡敌方攻击。防御手势本身不造成伤害，但会提供较高的防御力，使某些攻击无法生效。

防御手势本身也可能成为别人攻击的对象。
若某攻击打向防御手势，即使没有打穿，某些情况下仍然会被视为完成了一次攻击交互。这一部分会在后文的多人结算规则中详细说明。

\subsubsection{十字防}

手势为：双手握拳置于身前，左手小臂平行置于右手小臂上方。

消耗 2 格气，防御力为 3。

十字防是基础防御手势之一，能够有效抵挡许多中等强度的攻击。

\subsubsection{八卦}

手势为：双手小臂置于身前上下移动，手掌面向自己，手指自由弯曲。

消耗 3 格气，防御力为 4。

八卦是比十字防更高一级的防御手势，它能抵挡更多高强度攻击。

\subsection{锦囊手势}

锦囊手势属于特殊功能类手势。
它们不依靠普通的攻击力大于防御力来发挥作用，而是通过取消攻击、规避结算、获得镐等方式影响局势。
锦囊中有些手势可以针对特定攻击进行反制，有些则作用于自己。

\subsubsection{你吃}

手势为：伸出食指指向前方。

消耗 1 格气，防御力为 0，可以作用于特定攻击手势进行反制：

\begin{itemize}
    \item 闪电：若你吃成功作用于闪电，则闪电直接失效，且其使用者无法获得电池。
    \item 破：若你吃成功作用于破，则破直接失效，且破的使用者会受到 1 点反噬伤害。
\end{itemize}

\subsubsection{双吃}

手势为：伸出食指和中指并拢指向前方。

消耗 2 格气，防御力为 0，可以针对特定攻击进行反制：

\begin{itemize}
    \item 闪电：若双吃成功作用于闪电，则闪电直接失效，且其使用者无法获得电池。
    \item 破：若双吃成功作用于破，则破直接失效，且破的使用者会受到 1 点反噬伤害。
    \item Shining：若双吃成功作用于 Shining，则 Shining 的攻击直接失效。
\end{itemize}

双吃还可以拆分为两个你吃分别结算，作用对象变成两个分别判断的目标。
与你吃相同，双吃若未命中目标，则资源仍然消耗且玩家不会成为已操作对象。

\subsubsection{闪}

手势为：伸出一根手指在身前晃动。

无消耗，每局最多使用两次。

闪是一种完全规避型的锦囊。玩家使用闪后，会在结算正式开始前立即生效，并直接成为已操作对象。
使用闪的玩家在本回合中不能被任何攻击或锦囊选为目标，也不再参与本回合的其他结算。

\subsubsection{加镐}

手势为：右手握拳锤向心脏位置。

使用后获得一个镐，防御力为 0。

\subsection{特殊规则}

在了解了各个手势的基本内容后，我们还需要介绍几个特殊规则。

\subsubsection{gi攻击黑洞}

gi可以攻击黑洞，这是唯一一个低攻击力手势可以攻击高攻击力手势的情况。

当 gi 成功攻击黑洞时，黑洞使用者会遭到反噬并受到 3 点伤害。

\subsubsection{gi抢镐}

gi 可以对本回合出加镐的玩家发动抢镐。

抢镐成功后，gi 使用者与出镐者同时成为已操作对象，gi 使用者获得一个镐，而出镐者无法获得镐。

\subsubsection{gi的其它特殊规则}

gi 不能攻击出 gi 的玩家。也就是说，若某玩家本回合出的也是 gi，则另一名玩家不能选择他作为 gi 的攻击目标。

与其他攻击手势不同，若 gi 存在任何可攻击目标，则 gi 不能放空，必须选择攻击。
其他攻击手势通常可以根据局势选择放空，只保留自己的防御力；但 gi 只要能打，就必须打。
若 gi 没有任何可攻击目标，则该 gi 视为失效，且使用者保持未操作对象。

\subsubsection{黑洞、Shining与双吃的拆分}

黑洞可以拆分为三个小黑洞。
拆分后，每个小黑洞造成 1 点伤害，但攻击力仍视为 5。
三个小黑洞的目标必须在结算时一次性全部分配完毕，不能先打一段再根据结果决定剩下两段打谁。
多个小黑洞可以分配给同一个目标，因此黑洞可以一次性对同一名玩家造成 2 点甚至 3 点伤害。
由于黑洞拆分后依然属于黑洞结算，因此其速度不变，被黑洞命中的多个目标会同时成为已操作对象。

Shining 可以拆分为两个闪电分别结算。
拆分后，每一段都按闪电来处理，包括攻击力、可能被吃掉的关系以及结算速度。

双吃可以拆分为两个你吃分别结算。
拆成两个你吃后，本质上是进行两次独立的你吃结算。

%------------------------------------------------
\section{结算机制}

多人版与双人版最根本的区别，主要集中在结算阶段。
双人版中，许多情况天然只有一个对手，因此目标选择与结算顺序都比较直接；
多人版中，场上同时存在多个潜在目标、多个同速手势与多个可能互相影响的结算结果，因此必须先建立一套统一的结算框架。

\subsection{操作对象}

当所有玩家完成出手，正式进入结算阶段时，场上所有玩家都成为未操作对象。
若某名玩家在结算过程中已经完成了一次有效交互，则该玩家会变成已操作对象。

并不是所有出手动作都会自动让玩家成为已操作对象。只有真正发生了有效交互，相关玩家才会进入已操作对象。
通常来说，以下情况列为有效交互，会让玩家成为已操作对象：

\begin{itemize}
    \item 成功攻击了别人或被别人攻击。
    \item 使用锦囊成功作用于别人、被别人的锦囊成功作用或使用作用于自身的锦囊。
    \item 参与了三连。
\end{itemize}

成功攻击别人并不需要造成伤害，有些未造成伤害的情况，依然可以使玩家变成已操作对象。
例如，某攻击手势去攻击一个防御手势，虽然没有打穿，但攻击行为本身已经成立，此时双方仍会成为已操作对象。
又例如，两名玩家使用同攻击力的攻击手势互相攻击时，虽然最终伤害结果等价于双方都没有掉血，
但这次相互攻击本身已经成立，因此双方依然成为已操作对象。

与之相对，有些行为虽然已经出手，但如果没有形成有效交互，则不会使玩家进入已操作状态。例如：
你吃或双吃没有命中任何目标、某个攻击没有选中任何目标、gi 在本回合没有任何可攻击目标等。
这些情况会消耗对应的资源，但玩家仍保持未操作对象状态。

\subsection{三连}

三连是拍拍 2.0 多人版中最特殊的一类结构，是一种由三名玩家之间的可伤害关系构成的循环。
三连在闪结算完毕之后立即判定，仅影响参与这组三连的三个人，不影响场上其他玩家的继续结算。

若三名玩家 A、B、C 满足：A 能让 B 扣血、B 能让 C 扣血、C 能让 A 扣血，则称这三人形成一个三连。
这里“能让对方扣血”指的是按照当前回合各自出的手势与单纯克制关系能造成伤害，而不考虑其他因素。

三连主要分为两种类型：gi — 你吃/双吃 — 破 、gi — 黑洞 — 其它攻击，
这里的其它攻击指除 gi 与黑洞外的攻击手势。

当三连成立时，三人之间的攻击全部无效并立即成为已操作对象。
若三连中的某人出的是闪电，由于闪电被吃掉，则该闪电在三连中不能获得电池。
但若三连中的某人出的是 Fire，则该 Fire 仍可以获得火种。

\subsubsection{三连人选的处理}

并不是所有时候场上都恰好只有三个互相对应的人。
在多人局中，可能出现某一类位置上有一个人，某一类位置上有两个人甚至更多人的情况。
三连人数结构主要按以下方式处理：

\begin{itemize}
    \item 若三方人数均为 1，即形成（1，1，1），则这三个人直接构成三连。
    \item 若三方中只有一方人数大于等于 2，即形成（1，1，多），则由克制这一“多人方”的那一方，
    从多人方中选择一人，再与另外两方共同组成三连。
    \item 若三方中只有一方为 1 人，而另外两方都为多人，即形成（1，多，多），
    则先由这个“单人方”在自己所克制的所有对象中选择一人，
    然后由被选中的这人继续在自己所克制的所有对象中再选择一人，最终由这三人构成三连。
\end{itemize}

这种处理方式的目的是在多人局中把原本只有三人的循环关系推广到更复杂的分布情形下，同时保证最终真正进入三连结算的仍然只是三个人。

\subsubsection{出现两组三连}

若场上可以独立组成两组三连（共六名玩家），则本回合直接结束，除资源手势与加镐仍正常结算外，其余攻击和锦囊均无效。

\subsection{速度}

当结算阶段开始时，所有玩家出手的手势会被划分到不同的速度层中。
速度的快慢决定了结算的先后顺序，速度快的手势先结算，速度慢的手势后结算。
具体速度关系如下：

\begin{itemize}
    \item 闪
    \item 三连
    \item 你吃 / 双吃
    \item 攻击黑洞的 gi
    \item 黑洞
    \item 如来 / Shining
    \item 冷锋 / 烈焰
    \item 攻击其他目标造成伤害的 gi / 抢镐的 gi
    \item 破 / 闪电
    \item Fire
    \item 没有目标的 gi
    \item 气 / 盾 / 加镐
\end{itemize}

gi在不同状态下具有不同的速度。当结算进入某一 gi 所在的速度层时，需要先确认其本回合的攻击对象；
若其目标不属于该速度层对应的对象范围，则该 gi 不在此轮结算。

\subsection{同速下的结算机制}

若多个玩家位于同一速度层，同速玩家先要各自说明自己此时想操作的对象。
然后这些玩家之间可以进行协商与调整。这个协商过程可以不止一轮，
玩家可以根据他人意图反复修改自己的目标，直到达成一致的可执行方案。

如果同速玩家之间最终没有冲突，例如 A 想打 B，C 想打 D，而且这些目标互不重叠，
也不违反其他规则，那么这些手势就同时打出、同时结算，并进入下一个速度层。

\subsubsection{选择目标}

攻击与锦囊在选择目标时，必须同时满足以下原则：

\begin{itemize}
    \item 目标必须是未操作对象；
    \item 目标必须符合该手势本身允许作用的范围，例如你吃只能选择出闪电或者破的玩家作为目标、闪和加镐选择自己作为目标等。
    \item 目标必须符合其它规定，如gi的特殊规则、攻击手势不能攻击攻击力高于自己的攻击手势等。
\end{itemize}

\subsubsection{同速冲突}

同速冲突大致有两类，一类是一般冲突，另一类是多攻击少冲突。

一般冲突指若干同速玩家之间的目标安排彼此不兼容，但并非出现多人同时争抢攻击同一个目标的局面。
在这种情况下，玩家继续协商；若协商失败，其中一部分没有谈拢的玩家放空，而另外一部分已经谈拢的玩家仍按约定执行。

多攻击少冲突指的是在同一个速度层中，出现两名或以上玩家想攻击同一名玩家的情形。
此时由被攻击方决定“由谁来攻击自己”，若围绕这个分配仍无法协商出结果，则随机决定。

当黑洞、Shining和双吃拆分后，需要同时为每一段选择目标，
并且这些目标必须在同一轮协商中一次性全部确定。

\subsubsection{同速攻击与防御归零}

一旦玩家决定发动攻击，则其防御力会在攻击打出的同时变为 0。这条规则对同速层尤其重要。
例如，两个同速的攻击者 A 与 B 若都出的是闪电，那么如果 A 决定攻击某个低防御目标C，
而 B 决定攻击 A，那么 A 在把闪电向C打出去的同时，自己原本的防御就会归零，
这使得 B 的攻击有机会命中 A，这称作 A 强杀 C。

因此，同速玩家之间必须先交流与协商，决定最终如何安排攻击，不能简单理解成先选先得或各自独立结算。
若同速玩家之间最终无冲突，则这些攻击视为同时发出、同时命中、同时扣血。若有人因此死亡，也是在同一时刻同时死亡。

\subsubsection{放空规则}

除 gi 外的攻击手势都可以选择放空，即玩家本回合保留自己原本手势所提供的防御力，但不对任何其他人发动攻击。
你吃和双吃也可以放空，即不选择任何目标，但这意味着其防御力为 0，因为这两个手势本身的防御力就是 0。

gi 若有任何可攻击目标，则不能放空，必须从合法目标中选择一个进行攻击。
只有当 gi 不存在任何合法目标时，才视为失效，并保持未操作对象状态。

对于任何手势来说，如果没有合法目标，从结果上和放空相同。

\subsection{结算结束}

即使玩家出气、盾、加镐、Fire、闪电在轮到其速度级之前就已经成为了已操作对象，
其仍然能够获得相应的资源（被吃掉的闪电和被抢镐除外）。

当按速度结算完所有手势之后，生命值降至 0 或以下的玩家视为死亡，与“蛤蟆”“蟆蛤”“蚂蚁”一起从下回合开始移出游戏。

%------------------------------------------------
\section{总结}

\textbf{此段由ChatGPT提供}

拍拍（Clapclap）2.0 在基础手势系统的基础上，进一步补充了多人局中的完整结算规则。
游戏仍然以同步出手为核心，但在结算阶段中引入了速度、未操作对象、已操作对象与三连等机制，
使多人局中的目标选择、冲突处理与伤害结算能够按照统一规则进行。

在本规则体系下，玩家需要先理解各类手势本身的资源消耗、攻击力、防御力与特殊效果，再结合多人局中的速度顺序与目标合法性进行判断。
资源手势决定后续发展的基础，攻击手势提供主要伤害来源，防御手势用于抵挡高强度攻击，
锦囊手势则通过取消攻击、规避结算或改变资源状态来影响局势。

与双人局相比，多人局中更需要注意“谁先结算”“谁仍然是未操作对象”“同速玩家之间是否发生冲突”“三连是否已经成立”等问题。
许多看似相同的出手，在不同人数、不同速度层和不同已操作状态下，结算结果可能并不相同。
因此，在实际游戏过程中，先按规则确认速度层，再按未操作对象进行目标选择，最后处理伤害、资源与死亡，通常是最清晰的结算顺序。

当玩家熟悉了手势本身与多人版结算逻辑之后，就可以按照本规则完成完整对局。
若后续还要继续扩展版本，也应尽量保持现有术语、速度层次与目标选择规则的一致性，以便不同版本之间能够平稳衔接。

\newpage
\section*{修订记录}
\addcontentsline{toc}{section}{修订记录}

\begin{center}
\begin{tabular}{|c|c|c|c|}
\hline
版本号 & 日期 & 修订人 & 修订内容 \\
\hline
v1.0.0 & 2026年3月20日 & Yan & 初始版本完成。 \\
\hline
v1.0.1 & 2026年3月20日 & Yan & 修正了部分表述错误 \\
\hline
\end{tabular}
\end{center}

\end{document}